
\documentclass[12pt,a4paper]{article}

% ------------------------------------------------------------
% Fonts & layout
% ------------------------------------------------------------
\usepackage[T1]{fontenc}
\usepackage[utf8]{inputenc}
\usepackage{lmodern}
\usepackage{microtype}
\usepackage{geometry}
\geometry{margin=1in}
\usepackage{array}
\usepackage{booktabs}
\usepackage[table]{xcolor}
\usepackage{tikz}
\usetikzlibrary{positioning}

% ------------------------------------------------------------
% Math
% ------------------------------------------------------------
\usepackage{amsmath,amssymb,amsthm,mathtools}

% Standard operators
\DeclareMathOperator{\id}{id}
\DeclareMathOperator{\im}{im}
\DeclareMathOperator{\coker}{coker}
\DeclareMathOperator{\Hom}{Hom}

\DeclareMathOperator{\Spec}{Spec}
\DeclareMathOperator{\Proj}{Proj}

% Categories / notation
\newcommand{\Ab}{\mathbf{Ab}}
\newcommand{\PreSh}{\mathbf{PreSh}}
\newcommand{\Sh}{\mathbf{Sh}}

\newcommand{\F}{\mathcal F}


% basic number sets
\newcommand{\RR}{\mathbb{R}}

% sheaf letters

\newcommand{\G}{\mathcal G}

\newcommand{\GG}{\mathbb G}


\newcommand{\AffSch}{\mathbf{AffSch}}
\newcommand{\ZZ}{\mathbb Z}

\newcommand{\PP}{\mathbb P}

\usepackage{tikz}
\usetikzlibrary{decorations.pathreplacing}
\usetikzlibrary{positioning,arrows.meta}

% sheaf of continuous real-valued functions
\newcommand{\Czero}{\mathcal C^0}        % “C^0”
\newcommand{\CzeroX}{\mathcal C^0_X}     % sheaf on X (optional)

\newcommand{\Cinfty}{\mathcal C^\infty}
\newcommand{\CinftyM}{\mathcal C^\infty_M} % if you use C^\infty_M

% number sets / standard symbols
\newcommand{\CC}{\mathbb{C}}

% structure sheaf notation in complex geometry
\newcommand{\OO}{\mathcal{O}}
\newcommand{\OOX}{\mathcal{O}_X} % optional convenience



\newcommand{\Sch}{\mathrm{Sch}}
\newcommand{\Rings}{\mathrm{Rings}}

\newcommand{\kappares}{\kappa}
\newcommand{\GammaX}{\Gamma(X,\mathcal O_X)}

\newtheorem{theorem}{Theorem}[section]
\newtheorem{proposition}[theorem]{Proposition}
\newtheorem{lemma}[theorem]{Lemma}

\newcommand{\nilrad}{\sqrt{(0)}}

\newcommand{\Frac}{\operatorname{Frac}}
\newcommand{\codim}{\operatorname{codim}}
\newcommand{\trdeg}{\operatorname{trdeg}}

\newcommand{\height}{\operatorname{height}}

\newcommand{\Pic}{\operatorname{Pic}}
\newcommand{\divisor}{\operatorname{div}}

\newcommand{\Z}{\mathbb{Z}}
\newcommand{\Cl}{\operatorname{Cl}}

\providecommand{\Spec}{\operatorname{Spec}}
\providecommand{\A}{\mathbb A}
\providecommand{\Pj}{\mathbb P}
\providecommand{\G}{\mathbb G}

% ------------------------------------------------------------
% Links
% ------------------------------------------------------------
\usepackage{xcolor}
\usepackage{hyperref}
\usepackage{bookmark}
\hypersetup{
  colorlinks=true,
  linkcolor=blue!50!black,
  urlcolor=blue!50!black,
  citecolor=blue!50!black,
  pdftitle={Algebraic Geometry 20
  },
  pdfauthor={},
  pdfsubject={Lecture notes (reorganized)},
  pdfcreator={TeXStudio / LaTeX}
}

% ------------------------------------------------------------
% Diagrams
% ------------------------------------------------------------
\usepackage{tikz-cd}

% ------------------------------------------------------------
% Lists
% ------------------------------------------------------------
\usepackage{enumitem}
\setlist[itemize]{topsep=4pt,itemsep=2pt,parsep=0pt,partopsep=0pt}
\setlist[enumerate]{topsep=4pt,itemsep=2pt,parsep=0pt,partopsep=0pt}

% ------------------------------------------------------------
% Theorem environments
% ------------------------------------------------------------
\newtheoremstyle{notespace}{6pt}{6pt}{\normalfont}{}{\bfseries}{.}{0.5em}{}
\theoremstyle{notespace}
\newtheorem{definition}{Definition}[subsection]

\newtheorem{corollary}[definition]{Corollary}
\newtheorem{remark}[definition]{Remark}
\newtheorem{example}[definition]{Example}
\newtheorem{warning}[definition]{Warning}


\DeclareMathOperator{\Log}{Log}

\newcommand{\B}{\mathcal B}
\newcommand{\C}{\mathcal C}
\newcommand{\Ocal}{\mathcal O}
\newcommand{\I}{\mathcal I}


\newcommand{\Fp}{\mathcal F'}      % subsheaf
\newcommand{\Fpp}{\mathcal F''} 

\newcommand{\GammaU}[1]{\Gamma(U,#1)}

\DeclareMathOperator*{\colim}{colim}

% circle
\newcommand{\SSone}{S^1}


\usetikzlibrary{arrows.meta,positioning,calc,shapes.geometric,fit}

\usepackage{adjustbox}

% sheaves of continuous functions


% ------------------------------------------------------------
% Paragraph spacing
% ------------------------------------------------------------
\setlength{\parskip}{0.5em}
\setlength{\parindent}{0pt}

\setcounter{tocdepth}{2}
\setcounter{secnumdepth}{2}

% ------------------------------------------------------------
% Fancy headers
% ------------------------------------------------------------
\usepackage{fancyhdr}
\pagestyle{fancy}
\fancyhf{}
\lhead{Theory of Algebraic Geometry 20}
\rhead{\thepage}
\renewcommand{\headrulewidth}{0.4pt}
\setlength{\headheight}{14.5pt}




\begin{document}


\title{Algebraic Geometry 20\\
	\large Riemann surfaces}
\author{}
\date{}

	\maketitle
	

	\setcounter{section}{0}

\tableofcontents

	\clearpage

\section{Curvature Tensors: From Metric to Spacetime}

The main computational and conceptual pipeline is
\[
g
\longrightarrow
\Gamma
\longrightarrow
\text{geodesics}
\longrightarrow
Riemann
\longrightarrow
Ricci
\longrightarrow
\text{scalar curvature}
\longrightarrow
Weyl.
\]

The metric \(g\) is the starting object. Everything else is computed from it.

\subsection{Level 1: Metric Tensor}

A metric tensor tells us how to measure lengths, angles, and distances. In local coordinates
\[
x^1,\dots,x^n,
\]
it is written as
\[
ds^2=g_{ij}\,dx^i dx^j.
\]

For tangent vectors
\[
X=X^i\partial_i,
\qquad
Y=Y^j\partial_j,
\]
the metric gives
\[
g(X,Y)=g_{ij}X^iY^j.
\]

Thus the metric is a machine of the form
\[
\boxed{
	\text{two vectors}\longrightarrow \text{one number}.
}
\]

For example, if
\[
g=
\begin{pmatrix}
	1&0\\
	0&1
\end{pmatrix},
\qquad
X=(1,2),
\qquad
Y=(3,4),
\]
then
\[
g(X,Y)=1\cdot 3+2\cdot 4=11.
\]

So the metric is the first object needed before curvature can be discussed.

\subsection{Level 2: Christoffel Symbols}

The Christoffel symbols describe how coordinate directions change from point to point. They are computed from the metric by
\[
\Gamma^k_{ij}
=
\frac12 g^{k\ell}
\left(
\partial_i g_{j\ell}
+
\partial_j g_{i\ell}
-
\partial_\ell g_{ij}
\right).
\]

They are not tensors, but they are essential for computing geodesics and curvature.

Conceptually,
\[
\boxed{
	\text{metric changes}
	\longrightarrow
	\text{Christoffel symbols}.
}
\]

If the metric is constant, then all derivatives vanish, and hence
\[
\Gamma^k_{ij}=0.
\]

This is what happens in ordinary Euclidean coordinates.

\subsection{Level 3: Geodesics}

Geodesics are the straightest possible paths. They satisfy
\[
\frac{d^2x^k}{dt^2}
+
\Gamma^k_{ij}
\frac{dx^i}{dt}
\frac{dx^j}{dt}
=0.
\]

If
\[
\Gamma^k_{ij}=0,
\]
then
\[
\frac{d^2x^k}{dt^2}=0,
\]
so geodesics are straight lines.

If Christoffel symbols are nonzero, geodesics bend. Thus
\[
\boxed{
	\Gamma
	\longrightarrow
	\text{geodesic equations}.
}
\]

\[
\begin{array}{c|c}
	\text{Geometry} & \text{Geodesics}\\
	\hline
	\text{Flat plane} & \text{straight lines}\\
	\text{Sphere} & \text{great circles}\\
	\text{Hyperbolic plane} & \text{curved geodesics diverging from each other}\\
	\text{Spacetime} & \text{worldlines and light rays}
\end{array}
\]

\subsection{Level 4: Riemann Curvature Tensor}

The Riemann tensor measures true intrinsic curvature. It is computed from the Christoffel symbols:
\[
R^\rho_{\ \sigma\mu\nu}
=
\partial_\mu\Gamma^\rho_{\nu\sigma}
-
\partial_\nu\Gamma^\rho_{\mu\sigma}
+
\Gamma^\rho_{\mu\lambda}\Gamma^\lambda_{\nu\sigma}
-
\Gamma^\rho_{\nu\lambda}\Gamma^\lambda_{\mu\sigma}.
\]

It takes three vectors and returns a vector:
\[
R(X,Y)Z.
\]

Geometrically, it measures what happens when a vector is parallel transported around a tiny loop.

\[
\boxed{
	\text{three vectors}
	\longrightarrow
	\text{one vector}.
}
\]

After lowering an index, the tensor
\[
R_{ijkl}
\]
takes four vectors and returns one number.

\subsection{Level 5: Ricci Tensor}

The Ricci tensor is obtained by taking a trace of the Riemann tensor:
\[
Ric_{\sigma\nu}
=
R^\rho_{\ \sigma\rho\nu}.
\]

So Ricci is a compressed version of Riemann. It takes two vectors and returns a number:
\[
Ric(X,Y)=Ric_{ij}X^iY^j.
\]

Thus
\[
\boxed{
	\text{two vectors}
	\longrightarrow
	\text{one number}.
}
\]

Geometrically, Ricci curvature measures averaged curvature around a direction.

\[
\begin{array}{c|c}
	\text{Ricci sign} & \text{Geometric meaning}\\
	\hline
	Ric>0 & \text{nearby geodesics tend to focus}\\
	Ric<0 & \text{nearby geodesics tend to diverge}\\
	Ric=0 & \text{no Ricci volume focusing}
\end{array}
\]

A Ricci matrix is the table
\[
Ric_{ij}=Ric(e_i,e_j).
\]

The entries are not chosen arbitrarily; they are obtained by contracting the Riemann tensor.

\subsection{Level 6: Scalar Curvature}

The scalar curvature is the trace of the Ricci tensor:
\[
R=g^{ij}Ric_{ij}.
\]

It is one number at each point:
\[
\boxed{
	\text{point}
	\longrightarrow
	\text{one curvature number}.
}
\]

In two dimensions,
\[
R=2K,
\]
where \(K\) is Gaussian curvature.

\[
\begin{array}{c|c}
	\text{Geometry} & \text{Scalar curvature}\\
	\hline
	\text{sphere} & \text{positive}\\
	\text{plane} & 0\\
	\text{hyperbolic plane} & \text{negative}
\end{array}
\]

\subsection{Level 7: Weyl Tensor}

The Weyl tensor is the trace-free part of the Riemann tensor. For \(n\geq 3\),
\[
C_{ijkl}
=
R_{ijkl}
-
\frac{1}{n-2}
\left(
g_{ik}Ric_{jl}
-
g_{il}Ric_{jk}
-
g_{jk}Ric_{il}
+
g_{jl}Ric_{ik}
\right)
+
\frac{R}{(n-1)(n-2)}
\left(
g_{ik}g_{jl}
-
g_{il}g_{jk}
\right).
\]

Conceptually,
\[
\boxed{
	Riemann
	=
	Weyl
	+
	Ricci\ part
	+
	scalar\ part.
}
\]

Ricci measures volume focusing, while Weyl measures shape distortion.

\[
\begin{array}{c|c}
	\text{Curvature part} & \text{Effect on small ball}\\
	\hline
	\text{Ricci} & \text{small ball becomes smaller or larger}\\
	\text{Weyl} & \text{small ball becomes an ellipsoid with nearly same volume}
\end{array}
\]

Important facts:
\[
C=0
\]
in dimension \(2\), and also identically in dimension \(3\). Weyl curvature becomes genuinely important in dimension \(4\), especially in spacetime.

\subsection{Level 8: Spacetime Examples}

In spacetime, the metric has Lorentzian signature
\[
(-,+,+,+).
\]

Instead of ordinary distance, it measures spacetime interval:
\[
ds^2=g_{\mu\nu}dx^\mu dx^\nu.
\]

The same pipeline applies:
\[
g_{\mu\nu}
\longrightarrow
\Gamma^\rho_{\mu\nu}
\longrightarrow
R^\rho_{\ \sigma\mu\nu}
\longrightarrow
Ric_{\mu\nu}
\longrightarrow
R
\longrightarrow
C_{\rho\sigma\mu\nu}.
\]

In general relativity, Ricci curvature is tied to matter-energy through Einstein's equation, while Weyl curvature describes tidal and gravitational-wave type effects.

\section{Concrete Model Examples}

\subsection{Example A: Sphere of Radius \(r\)}

The sphere of radius \(r\) has metric
\[
ds^2=r^2d\theta^2+r^2\sin^2\theta\,d\phi^2.
\]

Thus
\[
g=
\begin{pmatrix}
	r^2&0\\
	0&r^2\sin^2\theta
\end{pmatrix},
\qquad
g^{-1}
=
\begin{pmatrix}
	\dfrac1{r^2}&0\\
	0&\dfrac1{r^2\sin^2\theta}
\end{pmatrix}.
\]

The nonzero Christoffel symbols are
\[
\Gamma^\theta_{\phi\phi}
=
-\sin\theta\cos\theta,
\]
and
\[
\Gamma^\phi_{\theta\phi}
=
\Gamma^\phi_{\phi\theta}
=
\cot\theta.
\]

The geodesic equations are
\[
\theta''
-
\sin\theta\cos\theta(\phi')^2
=0,
\]
and
\[
\phi''
+
2\cot\theta\,\theta'\phi'
=0.
\]

These describe great circles.

The Gaussian curvature is
\[
K=\frac1{r^2}.
\]

In two dimensions,
\[
Ric=Kg.
\]

Therefore
\[
Ric=\frac1{r^2}g.
\]

So
\[
Ric_{\theta\theta}=1,
\qquad
Ric_{\phi\phi}=\sin^2\theta.
\]

The scalar curvature is
\[
R=2K=\frac2{r^2}.
\]

The Weyl tensor is
\[
C=0,
\]
because in dimension \(2\), Weyl curvature is identically zero.

Hence
\[
\boxed{
	K=\frac1{r^2},
	\qquad
	Ric=\frac1{r^2}g,
	\qquad
	R=\frac2{r^2},
	\qquad
	C=0.
}
\]

For \(r=1\),
\[
K=1,
\qquad
R=2.
\]

For \(r=2\),
\[
K=\frac14,
\qquad
R=\frac12.
\]

A larger sphere is less curved.

\[
\text{Sphere radius }r
\longrightarrow
\text{color by }K=\frac1{r^2}
\longrightarrow
\text{great-circle geodesics}
\longrightarrow
\text{parallel transport rotation}.
\]

\subsection{Example B: Flat Plane}

The Euclidean plane has metric
\[
ds^2=dx^2+dy^2.
\]

Therefore
\[
g=
\begin{pmatrix}
	1&0\\
	0&1
\end{pmatrix}.
\]

Since \(g\) is constant,
\[
\Gamma^k_{ij}=0.
\]

Therefore
\[
R^i_{\ jkl}=0.
\]

Hence
\[
Ric=0,
\qquad
R=0,
\qquad
C=0.
\]

So
\[
\boxed{
	\Gamma=0,
	\qquad
	Riemann=0,
	\qquad
	Ric=0,
	\qquad
	R=0.
}
\]

Visualization:
\[
\text{straight geodesics}
\longrightarrow
\text{no parallel transport rotation}
\longrightarrow
\text{zero curvature color}.
\]

The plane is locally and globally flat.

\subsection{Example C: Hyperbolic Plane}

Use the upper half-plane model:
\[
ds^2=\frac{dx^2+dy^2}{y^2},
\qquad y>0.
\]

Then
\[
g=
\begin{pmatrix}
	\dfrac1{y^2}&0\\
	0&\dfrac1{y^2}
\end{pmatrix},
\qquad
g^{-1}
=
\begin{pmatrix}
	y^2&0\\
	0&y^2
\end{pmatrix}.
\]

The Gaussian curvature is
\[
K=-1.
\]

In dimension \(2\),
\[
Ric=Kg.
\]

So
\[
Ric=-g.
\]

Thus
\[
Ric=
\begin{pmatrix}
	-\dfrac1{y^2}&0\\
	0&-\dfrac1{y^2}
\end{pmatrix}.
\]

The scalar curvature is
\[
R=2K=-2.
\]

The Weyl tensor is again
\[
C=0,
\]
because it is a two-dimensional geometry.

Hence
\[
\boxed{
	K=-1,
	\qquad
	Ric=-g,
	\qquad
	R=-2,
	\qquad
	C=0.
}
\]

Visualization:
\[
\text{negative curvature}
\longrightarrow
\text{geodesics diverge}
\longrightarrow
\text{triangles have angle sum }<\pi
\longrightarrow
\text{constant negative curvature color}.
\]

Hyperbolic geometry is the standard model of constant negative curvature.

\subsection{Example D: Schwarzschild Spacetime}

The Schwarzschild metric is
\[
ds^2
=
-\left(1-\frac{2M}{r}\right)dt^2
+
\left(1-\frac{2M}{r}\right)^{-1}dr^2
+
r^2d\theta^2
+
r^2\sin^2\theta\,d\phi^2.
\]

Let
\[
f(r)=1-\frac{2M}{r}.
\]

Then
\[
g_{\mu\nu}
=
\begin{pmatrix}
	-f&0&0&0\\
	0&f^{-1}&0&0\\
	0&0&r^2&0\\
	0&0&0&r^2\sin^2\theta
\end{pmatrix}.
\]

Outside the mass, Schwarzschild is a vacuum solution:
\[
Ric_{\mu\nu}=0.
\]

Therefore
\[
R=0.
\]

But the Riemann tensor is not zero:
\[
R_{\rho\sigma\mu\nu}\neq0.
\]

Also,
\[
C_{\rho\sigma\mu\nu}\neq0.
\]

In fact, in vacuum Schwarzschild spacetime,
\[
C_{\rho\sigma\mu\nu}=R_{\rho\sigma\mu\nu}.
\]

Thus the curvature is pure Weyl curvature.

\[
\boxed{
	\text{Ricci curvature }=0
	\qquad
	\text{but}
	\qquad
	\text{tidal Weyl curvature }\neq0.
}
\]

This is important:
\[
\text{there is no matter locally outside the black hole, but spacetime is still curved.}
\]

A useful scalar invariant is the Kretschmann scalar:
\[
K_{\mathrm{Kr}}
=
R_{\alpha\beta\gamma\delta}
R^{\alpha\beta\gamma\delta}
=
\frac{48M^2}{r^6}.
\]

Hence
\[
\boxed{
	Ric_{\mu\nu}=0,
	\qquad
	R=0,
	\qquad
	Riemann\neq0,
	\qquad
	Weyl\neq0.
}
\]

Visualization:
\[
\text{black hole spacetime}
\longrightarrow
\text{Ricci color }=0
\longrightarrow
\text{Weyl/tidal intensity nonzero}
\longrightarrow
\text{light cones tilt}
\longrightarrow
\text{radial stretching and tangential compression}.
\]

Near the black hole, a small freely falling sphere becomes stretched radially and squeezed tangentially.

\subsection{Example E: FLRW Universe}

The FLRW metric is
\[
ds^2
=
-dt^2
+
a(t)^2
\left(
\frac{dr^2}{1-kr^2}
+
r^2d\theta^2
+
r^2\sin^2\theta\,d\phi^2
\right).
\]

Here
\[
a(t)
\]
is the scale factor, and
\[
k=1,0,-1
\]
represents positive, flat, or negative spatial curvature.

An important Ricci component is
\[
Ric_{00}
=
-3\frac{\ddot a}{a}.
\]

The scalar curvature is
\[
R
=
6\left(
\frac{\ddot a}{a}
+
\frac{\dot a^2+k}{a^2}
\right).
\]

Interpretation:
\[
a(t)
\longrightarrow
\text{space expands or contracts}
\longrightarrow
\text{Ricci curvature changes}
\longrightarrow
\text{geodesic distances change}.
\]

Unlike Schwarzschild vacuum, FLRW generally has nonzero Ricci curvature:
\[
\boxed{
	Ric\neq0
}
\]
in general, because the universe contains matter, radiation, dark matter, and dark energy in cosmological models.

Visualization:
\[
\text{expanding grid}
\longrightarrow
\text{scale factor }a(t)
\longrightarrow
\text{Ricci scalar heatmap over time}
\longrightarrow
\text{geodesic separation}
\longrightarrow
\text{cosmic expansion}.
\]

\section{Comparison Table}

\[
\begin{array}{c|c|c|c|c|c}
	\text{Example}
	&
	\text{Geometry}
	&
	K
	&
	Ric
	&
	R
	&
	C
	\\
	\hline
	\text{Flat plane}
	&
	\text{Euclidean surface}
	&
	0
	&
	0
	&
	0
	&
	0
	\\
	\text{Sphere }S^2_r
	&
	\text{positive curvature surface}
	&
	\dfrac1{r^2}
	&
	\dfrac1{r^2}g
	&
	\dfrac2{r^2}
	&
	0
	\\
	\text{Hyperbolic plane}
	&
	\text{negative curvature surface}
	&
	-1
	&
	-g
	&
	-2
	&
	0
	\\
	\text{Schwarzschild}
	&
	\text{vacuum spacetime}
	&
	\text{not surface }K
	&
	0
	&
	0
	&
	\neq0
	\\
	\text{FLRW universe}
	&
	\text{cosmology}
	&
	\text{depends on }k,a(t)
	&
	\neq0\ \text{usually}
	&
	\neq0\ \text{usually}
	&
	0\ \text{in ideal FLRW}
\end{array}
\]

\section{Visualization Levels for Math3D}

\subsection{Level 1: Metric Visualization}

Show local rulers. At every point, the metric determines a small ellipse or ellipsoid.

\[
\text{circle in coordinate space}
\longrightarrow
\text{metric transforms it}
\longrightarrow
\text{ellipse in measured geometry}.
\]

For the sphere, the metric changes with \(\theta\):
\[
g_{\phi\phi}=r^2\sin^2\theta.
\]

Near the pole, \(\phi\)-movement is small. Near the equator, \(\phi\)-movement is larger.

\subsection{Level 2: Christoffel Visualization}

Show basis vectors changing over the surface.

\[
\text{coordinate grid}
\longrightarrow
\text{basis vectors vary}
\longrightarrow
\text{Christoffel symbols appear}.
\]

For the sphere, lines of longitude meet at the poles, so the coordinate frame changes. That produces nonzero Christoffel symbols.

\subsection{Level 3: Geodesic Visualization}

Show geodesic curves.

\[
\begin{array}{c|c}
	\text{Geometry} & \text{Visualization}\\
	\hline
	\text{Plane} & \text{straight lines}\\
	\text{Sphere} & \text{great circles}\\
	\text{Hyperbolic plane} & \text{semicircles and vertical lines in upper half-plane}\\
	\text{Schwarzschild} & \text{light bending and orbital paths}
\end{array}
\]

\subsection{Level 4: Riemann Visualization}

Show parallel transport around a small loop.

\[
\text{start with arrow }Z
\longrightarrow
\text{move around tiny loop spanned by }X,Y
\longrightarrow
\text{compare final arrow with initial arrow}.
\]

The difference is curvature.

On a sphere, an arrow transported around a triangle rotates. On a plane, it does not rotate.

\subsection{Level 5: Ricci Visualization}

Show geodesic focusing.

\[
\begin{array}{c|c}
	\text{Ricci curvature} & \text{Visualization}\\
	\hline
	\text{positive Ricci} & \text{geodesics converge}\\
	\text{negative Ricci} & \text{geodesics diverge}\\
	\text{zero Ricci} & \text{no Ricci focusing}
\end{array}
\]

For a sphere, geodesics tend to meet. For hyperbolic space, geodesics tend to separate. For Schwarzschild vacuum, Ricci is zero outside the mass, but tidal distortion still exists.

\subsection{Level 6: Scalar Curvature Visualization}

Color the surface or spacetime slice by scalar curvature.

\[
\begin{array}{c|c}
	\text{Geometry} & \text{Scalar curvature}\\
	\hline
	\text{Sphere} & R=\dfrac2{r^2}>0\\
	\text{Plane} & R=0\\
	\text{Hyperbolic plane} & R=-2\\
	\text{Schwarzschild vacuum} & R=0\\
	\text{FLRW} &
	R=
	6\left(
	\dfrac{\ddot a}{a}
	+
	\dfrac{\dot a^2+k}{a^2}
	\right)
\end{array}
\]

\subsection{Level 7: Weyl Visualization}

Show tidal distortion.

\[
\text{small sphere of test particles}
\longrightarrow
\text{Weyl curvature}
\longrightarrow
\text{ellipsoid}.
\]

For Schwarzschild:
\[
\text{radial stretching}
\qquad
\text{and}
\qquad
\text{tangential compression}.
\]

This shows black-hole tidal forces. In two-dimensional surfaces, Weyl is zero, so this level becomes important mainly in spacetime.

\subsection{Level 8: Spacetime Visualization}

Display:
\[
\text{light cones},
\qquad
\text{worldlines},
\qquad
\text{null geodesics},
\qquad
\text{timelike geodesics},
\qquad
\text{tidal ellipsoids},
\qquad
\text{curvature invariants}.
\]

For Schwarzschild,
\[
K_{\mathrm{Kr}}=\frac{48M^2}{r^6}.
\]

So one can color space by tidal intensity.

For FLRW, the scale factor \(a(t)\) controls the expansion of space.

\section{Final Conceptual Summary}

The hierarchy is
\[
\boxed{
	g_{ij}
}
\]
which measures lengths,

\[
\boxed{
	\Gamma^k_{ij}
}
\]
which describes how coordinates and basis vectors change,

\[
\boxed{
	\text{geodesics}
}
\]
which are straightest paths determined by \(\Gamma\),

\[
\boxed{
	R^\rho_{\ \sigma\mu\nu}
}
\]
which measures full curvature,

\[
\boxed{
	Ric_{\mu\nu}
}
\]
which measures averaged curvature and geodesic focusing,

\[
\boxed{
	R
}
\]
which is one scalar curvature number,

and
\[
\boxed{
	C_{\rho\sigma\mu\nu}
}
\]
which measures trace-free tidal distortion.

For visualization:
\[
\begin{array}{c|c}
	\text{Object} & \text{Picture}\\
	\hline
	\text{Metric} & \text{local ruler}\\
	\text{Christoffel} & \text{changing frame}\\
	\text{Geodesics} & \text{straightest paths}\\
	\text{Riemann} & \text{loop rotation}\\
	\text{Ricci} & \text{focusing/diverging}\\
	\text{Scalar} & \text{curvature heatmap}\\
	\text{Weyl} & \text{tidal ellipsoid}\\
	\text{Spacetime} & \text{light cones and worldlines}
\end{array}
\]

\end{document}

